\section{Problem, scope, and contribution}
Let $K$ be a finite field, $D_i\subseteq K$ finite constrained domains, and
$f_1,\ldots,f_s\in K[X_1,\ldots,X_m]$. The generalized root-finding task is to
return the entire ordered set
\[
 {\cal Z}=\{(x_1,\ldots,x_m)\in D_1\times\cdots\times D_m:
                    f_j(x_1,\ldots,x_m)=0\text{ for every }j\}.
\]
An exact all-roots implementation must establish both membership and coverage.
If auxiliary variables occur in the input system, the requested projection and
whether auxiliary fibers must be returned are part of the input contract.
Finding a witness, completing a relaxation, and enumerating a projected norm
chart are different tasks.

The motivating application is the constrained polynomial-system stage of
elliptic-curve index calculus. Semaev's summation polynomials provide an algebraic
connection to point relations \cite{semaev}. The implementation considered here
instead searches a structured norm identity before reconstructing original
chain solutions. This manuscript documents that narrower implemented stage.
It is distinct from the repository's earlier \texttt{rootfinding-decomposition}
manuscript about direct subtraction and pair tables at smaller decomposition
orders. Its measurements and reported ratios are not imported here.

There are three concrete contributions. First, Theorem~\ref{thm:chart} gives
an exact binary norm-chart contract and a short coefficient proof that avoids
generic-position divisions. Second, Theorem~\ref{thm:compiler} explains why
bounded affine substitution changes representation without changing primary
models. Third, the measured native-coordinate profile and the subsequent
negative search result demonstrate that fewer declared variables need not
produce a hard-instance solution. These are mathematical exposition and
engineering results assembled from versioned artifacts, with no priority claim
for the underlying algebra or SAT techniques.

\subsection{Related work and the novelty boundary}
F4 organizes Gr\"obner-basis reductions using linear algebra
\cite{f4}; F5 introduces signature-based avoidance of reductions to zero
\cite{f5}. A comparison must identify an actual implementation, field backend,
input ideal, and final all-roots output, rather than compare algorithm labels.
The program also requires the repository-specific F6 baseline. This draft
does not identify a canonical F6 publication or establish that its available
interface matches the full root contract.

SAT with XOR reasoning and coordinated changes to representation are established
techniques \cite{cms}. Particularly close prior work is
Trimoska--Ionica--Dequen's WDSat approach to index-calculus systems
\cite{wdsat}, including ANF-aware parity reasoning and branching intended to
expose linear structure. Consequently, neither using SAT instead of Gr\"obner
bases nor branching on primary variables is a novelty claim here.
A full comparison of their mathematical model, preprocessing, and branching
against this norm parametrization remains necessary before submission.

Finite-field factorization and Frobenius computation are classical
\cite{cz,frobenius}. The support criterion below is an elementary consequence
of multiplicities in a quotient ring. Factoring an already supplied quintic
does not solve the search over its unknown coefficients. Our candidate research
advance would have to be a demonstrably useful structural reduction or pruning
rule beyond these ingredients. That advance, and its qualified end-to-end
performance, remain open.
